Yes, it does! The RLC formulas completely cover RL, RC, and LC circuits as well.
You do not need separate formulas for them. To find the formula for an RL, RC, or LC circuit, you simply take the comprehensive RLC formulas and set the missing component's value to zero (or remove it entirely).
Here is how you easily adapt the RLC equations for any of those subsets:
------------------------------
## ➡️ 1. How to Convert RLC to RL, RC, or LC## 📌 RL Circuits (No Capacitor)

* What to do: Remove the capacitive reactance term completely ($X_C = 0$).
* Series Impedance: From $Z = \sqrt{R^2 + (X_L - X_C)^2}$ ➡️ drops down to $Z = \sqrt{R^2 + X_L^2}$.
* Phase Angle: $\theta = \tan^{-1}\left(\frac{X_L}{R}\right)$.

## 📌 RC Circuits (No Inductor)

* What to do: Remove the inductive reactance term completely ($X_L = 0$).
* Series Impedance: From $Z = \sqrt{R^2 + (X_L - X_C)^2}$ ➡️ drops down to $Z = \sqrt{R^2 + (-X_C)^2} = \sqrt{R^2 + X_C^2}$.
* Phase Angle: $\theta = \tan^{-1}\left(\frac{-X_C}{R}\right)$ (The negative sign indicates the current leads the voltage).

## 📌 LC Circuits (No Resistor / Ideal Resonance)

* What to do: Set the resistance to zero ($R = 0$).
* Series Impedance: From $Z = \sqrt{R^2 + (X_L - X_C)^2}$ ➡️ drops down to $Z = X_L - X_C$.
* At Resonance ($X_L = X_C$): Impedance $Z = 0$ Ω. A series LC circuit acts as a perfect short circuit at its resonant frequency!

------------------------------
## ➡️ 2. Summary Table of Series Sub-Circuits
If you ever want to see them isolated side-by-side, here is how the primary formulas look for each specific type:

| Circuit Type | Impedance Magnitude ($Z$) | Phase Angle ($\theta$) | Resonant Frequency ($f_0$) |
|---|---|---|---|
| RL | $\sqrt{R^2 + (\omega L)^2}$ | $\tan^{-1}\left(\frac{\omega L}{R}\right)$ | None (No oscillation) |
| RC | $\sqrt{R^2 + \left(\frac{1}{\omega C}\right)^2}$ | $\tan^{-1}\left(\frac{-1}{\omega C R}\right)$ | None (No oscillation) |
| LC | $\omega L - \frac{1}{\omega C}$ | $+90^\circ$ or $-90^\circ$ | $f_0 = \frac{1}{2\pi\sqrt{LC}}$ |
| RLC | $\sqrt{R^2 + \left(\omega L - \frac{1}{\omega C}\right)^2}$ | $\tan^{-1}\left(\frac{\omega L - \frac{1}{\omega C}}{R}\right)$ | $f_0 = \frac{1}{2\pi\sqrt{LC}}$ |

Are you currently studying these for an exam, or are you working on a specific lab/design problem? If you have actual values (like $R = 100\,\Omega$, $L = 10\,\text{mH}$), I can run the numbers through Python to show you exactly how these reductions work in practice.

